\section{TD d'application : mode A6, mise en conditions initiales}
\label{secTdA6}

\begin{UPSTIinfor}{Objectif}
Mettre en œuvre le mode GEMMA \textbf{A6} du robot, en mobilisant la
chaîne de conversion présentée ci-dessus. Les trous proposés sont de
taille croissante : les premiers portent sur les blocs de
\textbf{conversion} déjà présentés en détail, les derniers vous demandent
de \textbf{choisir vous-même} le mouvement le plus adapté.
\end{UPSTIinfor}

\subsection{Rappel du comportement attendu (mode A6)}

Le robot doit rejoindre sa position de référence au-dessus du guichet,
sans jamais entraver la circulation des palettes sur le convoyeur :

\begin{enumerate}
    \item monter l'axe \textbf{J3} (vertical) jusqu'à $Z=150\;\text{mm}$,
    sans modifier la position courante des autres axes ;
    \item rejoindre ensuite la position cartésienne complète
    $(X=90,\ Y=-410,\ Z=150,\ RX=-180,\ RY=0,\ RZ=0)$.
\end{enumerate}

Si un arrêt est demandé par l'opérateur pendant ce mode, le robot doit
être immobilisé (sans perdre la puissance) puis reprendre sa course dès
que la demande d'arrêt n'est plus active. Ce bouton est accessible dans
l'objet \texttt{hmiScara} via
\texttt{hmiScara.m\_Inputs.m\_nixStopRequest} : comme rappelé en
section~\ref{secConversionHmi}, ce contact est à ouverture (\emph{NC}),
donc un arrêt est \textbf{demandé} lorsque ce booléen vaut
\texttt{FALSE}.

\subsection{Blocs utilisés dans ce TD}

Les blocs \texttt{MC\_} de contrôle et de commande de mouvement partagent
tous le même objet \texttt{AxesGroup} (\texttt{T\_StaeubliRobot}, déjà
lu/écrit par \texttt{VAL\_ReadAxesGroup}/\texttt{VAL\_WriteAxesGroup}) et
un squelette d'entrées/sorties commun :

\begin{center}
\begin{tabular}{|l|l|p{7.5cm}|}
\hline
\textbf{Paramètre} & \textbf{Sens / type} & \textbf{Rôle} \\
\hline
\texttt{Execute} & \texttt{VAR\_INPUT}, \texttt{BOOL} & Déclenche
l'exécution sur \textbf{front montant}. \\
\hline
\texttt{AxesGroup} & \texttt{VAR\_IN\_OUT}, \texttt{T\_StaeubliRobot} &
Objet du robot, commun à tous les blocs \texttt{MC\_}. \\
\hline
\texttt{Busy} / \texttt{Done} / \texttt{Error} & \texttt{VAR\_OUTPUT},
\texttt{BOOL} & En cours d'exécution / terminé avec succès / terminé en
erreur. \\
\hline
\end{tabular}
\end{center}

Les blocs de \textbf{mouvement} (\texttt{MC\_MoveLinearAbsolute},
\texttt{MC\_MoveDirectAbsolute}) ajoutent notamment, en \texttt{VAR\_INPUT} :
\texttt{Position} (\texttt{T\_CartesianPos}, la cible), \texttt{Velocity},
\texttt{Acceleration}, \texttt{Deceleration} (en~\% de la valeur
nominale), \texttt{CartesianVel}/\texttt{RotationalVel} (vitesses max. du
TCP) et \texttt{CoordSystem}/\texttt{ToolNumber} (repère et outil utilisés
-- \texttt{0} désigne respectivement le repère \texttt{world} et l'outil
\texttt{flange}).

\texttt{MC\_GroupPower} n'a qu'une entrée \texttt{Enable} (pas
\texttt{Execute}) et une seule sortie \texttt{Status} (état effectif de la
puissance) en plus d'\texttt{Error}/\texttt{ErrorID} -- il n'a pas de
notion de front, le bras reste sous/hors puissance tant que
\texttt{Enable} est à \texttt{TRUE}/\texttt{FALSE}.

\texttt{MC\_GroupInterrupt} et \texttt{MC\_GroupContinue} suivent le
squelette commun (\texttt{Execute}, \texttt{AxesGroup}, \texttt{Busy},
\texttt{Done}, \texttt{Error}) : le premier arrête le mouvement en cours
\emph{sans} vider la pile de commandes, le second relance l'exécution de
cette pile.

\UPSTIremarque[Objet \texttt{TCtrlStatusGroup}]{Par souci de lisibilité,
les sorties \texttt{Busy}/\texttt{Done}/\texttt{Error}/\texttt{ErrorID} de
ces blocs sont habituellement regroupées, lors de l'appel, dans un objet
de type \texttt{TCtrlStatusGroup} (champs \texttt{m\_xEnable},
\texttt{m\_xBusy}, \texttt{m\_xDone}, \texttt{m\_xError},
\texttt{m\_udiErrorID}, \texttt{m\_xTimeOut}) -- même logique que pour
\texttt{THmi} en section~\ref{secConversionHmi} : un seul objet nommé
plutôt que de nombreuses variables éparpillées.}

\subsection{Niveau 1 -- Compléter les blocs de conversion}

\begin{UPSTIactivite}[4][Appel de \texttt{VAL\_ReadAxesGroup} et \texttt{VAL\_WriteAxesGroup}][][][TD]
\UPSTIquestion{Compléter les paramètres manquants (repérés par
\texttt{(* A\_COMPLETER *)}) dans la section \texttt{ReadEthIpRobot} et
dans la section \texttt{WriteEthIpRobot} ci-dessous.}

\begin{lstlisting}
(* Section ReadEthIpRobot -- premiere section executee dans le cycle *)
InstReadAxesGroup(
    Enable     := (* A_COMPLETER *),
    AxesGroup  := (* A_COMPLETER *),
    IN         := fromRobotScara
);

(* ... section WriteEthIpRobot -- derniere section executee dans le cycle *)
InstWriteAxesGroup(
    Enable     := NOC_ETHIP_ROBOT_IN.HEALTH_BITS_IN[0].0,
    AxesGroup  := staeubliRobotScara,
    OUT        => (* A_COMPLETER *)
);
\end{lstlisting}
\vspace{2cm}
\end{UPSTIactivite}

\subsection{Niveau 2 -- Compléter la mise sous puissance}

\begin{UPSTIactivite}[4][Mise sous puissance du robot][][][TD]
\UPSTIquestion{Le robot doit rester sous puissance pendant toute la durée
du mode A6. Compléter l'appel au bloc \texttt{MC\_GroupPower}
ci-dessous : l'entrée \texttt{Enable} doit être activée en permanence
dans ce mode, et l'objet \texttt{AxesGroup} doit être celui de la cellule.}

\begin{lstlisting}
InstGroupPowerA6(
    Enable     := (* A_COMPLETER : toujours actif en mode A6 *),
    AxesGroup  := (* A_COMPLETER *)
);
\end{lstlisting}
\vspace{2cm}
\end{UPSTIactivite}

\subsection{Niveau 3 -- Choisir et écrire les mouvements}

À ce stade, vous disposez de tous les éléments pour écrire vous-même la
séquence de mouvement. Deux blocs \texttt{MC\_} vus en TP peuvent
convenir : \texttt{MC\_MoveLinearAbsolute} (mouvement linéaire) et
\texttt{MC\_MoveDirectAbsolute} (interpolation articulaire vers une
position cartésienne cible).

\begin{UPSTIactivite}[4][Étape 1 -- monter l'axe J3][][][TD]
\UPSTIquestion{L'énoncé demande de monter l'axe \textbf{J3} à
$Z=150\;\text{mm}$ \textbf{sans modifier la position courante des autres
axes}. Sur un robot Scara, l'axe J3 est l'axe vertical : une cible
cartésienne qui ne change que $Z$, en conservant $X$, $Y$, $RX$, $RY$,
$RZ$ courants, produit exactement ce déplacement.

\begin{enumerate}
    \item Entre \texttt{MC\_MoveLinearAbsolute} et
    \texttt{MC\_MoveDirectAbsolute}, quel bloc choisissez-vous pour cette
    étape ? Justifier à partir de la trajectoire que doit suivre l'outil
    (une translation verticale pure, sans contrainte angulaire
    particulière).
    \item Écrire, en complétant le squelette ci-dessous, l'instanciation
    et l'appel de ce bloc. La position cible utilisera les coordonnées
    \textbf{courantes} du robot pour $X$, $Y$, $RX$, $RY$, $RZ$ (déjà
    disponibles dans \texttt{staeubliRobotScara.Status.CartesianPos}), et
    $Z=150$.
\end{enumerate}}

\begin{lstlisting}
targetPosA6Etape1.X  := staeubliRobotScara.Status.CartesianPos.X;
targetPosA6Etape1.Y  := (* A_COMPLETER *);
targetPosA6Etape1.Z  := (* A_COMPLETER : valeur numerique de l'enonce *);
targetPosA6Etape1.RX := (* A_COMPLETER *);
targetPosA6Etape1.RY := (* A_COMPLETER *);
targetPosA6Etape1.RZ := (* A_COMPLETER *);

(* A_COMPLETER : instanciation et appel du bloc MC_ choisi,
   avec Execute, AxesGroup et la position cible ci-dessus *)
\end{lstlisting}
\vspace{3cm}
\end{UPSTIactivite}

\begin{UPSTIactivite}[4][Étape 2 -- rejoindre la position de référence][][][TD]
\UPSTIquestion{La position de référence
$(X=90,\ Y=-410,\ Z=150,\ RX=-180,\ RY=0,\ RZ=0)$ doit être rejointe par
une commande de mouvement \textbf{articulaire} (imposée par l'énoncé du
mode A6, contrairement à l'étape~1 où le choix était à justifier). Écrire
entièrement l'appel du bloc \texttt{MC\_} correspondant. Vous déclarerez
vous-même l'instance du bloc et la variable de position cible.}
\vspace{6cm}
\end{UPSTIactivite}

\begin{UPSTIactivite}[4][Gestion de l'arrêt demandé][][][TD]
\UPSTIquestion{Parmi les blocs \texttt{MC\_GroupInterrupt} et
\texttt{MC\_GroupContinue} vus en TP, lequel utiliser pour immobiliser le
robot sans perdre la puissance lorsqu'un arrêt est demandé, et lequel
pour reprendre la course lorsqu'il ne l'est plus ? Rédiger les deux
conditions d'appel (sur l'entrée \texttt{Execute} de chaque bloc, sans
écrire l'instanciation complète), en utilisant
\texttt{hmiScara.m\_Inputs.m\_nixStopRequest}.}
\vspace{4cm}
\end{UPSTIactivite}
